% lua-tikz3dtools-doc-problemstatement
\chapter{Problem Statement}

\section{Background and Context}
% Write 2–3 paragraphs that:
% - Introduce the general field or domain (e.g., mathematics education, visualization systems, applied linear algebra).
% - Summarize the state of practice or research that motivates this work.
% - Provide just enough context so a reader outside the field understands why this area is worth studying.

The field to which \luatikztdtools{} belongs is inherently
multidisciplinary, combining elements of programming and
mathematics, especially linear algebra. In my view, the current
state of practice remains fragmented. Much of the existing
literature focuses on isolated components of the broader problem;
for example, publications may address tessellation of surfaces or
the clipping of line segments by triangles, but seldom attempt a
unified framework. I believe I am the first to orchestrate linear
algebra and programming into a coherent synthesis that enables
parametric object generation tessellated by affine simplices
(points, line segments, and triangles). These simplices are then
systematically clipped against one another using a novel algorithm
that eliminates intersections. A second novel algorithm, also of
my design, performs occlusion sorting of the resulting simplices
via a transitive partial order.

This represents a genuinely groundbreaking achievement, with
applications in computer graphics, computational geometry, and
mathematical illustration.

\section{The Core Problem Being Addressed}
% Write 1 concise paragraph that:
% - States clearly and directly the central problem.
% - Avoids solutions or methods — only define what is problematic or missing.
% - Use plain, unambiguous wording so the reader knows exactly what challenge you are tackling.

The central problem addressed by \luatikztdtools{} is the
illustration of \(3\)D scenes composed of \(0\)--\(2\)-dimensional
affine simplices. Specifically, the package focuses on the
representation of parametric objects tessellated into points,
line segments, and triangles, including those that intersect.
The primary outstanding challenge is the elimination of cyclic
overlaps, which will be addressed in the next version of the
software.

\section{Scope and Boundaries of the Problem}
% Write 1–2 paragraphs that:
% - Define the specific scope (what is included in your investigation).
% - Explicitly mention what is *not* covered (to avoid unrealistic expectations).
% - Clarify assumptions you are making (e.g., focusing on 2D visualizations only, not 3D).

The scope of the problem---excluding cyclic overlaps---is twofold:
clipping intersecting simplices and partially transitive sorting
of the resulting non-intersecting simplices. This version of the
package assumes that cyclic overlaps do not occur. It is designed
for \(3\)D diagrams composed of tessellated parametric objects,
favoring exact partitioning over alternative methods to handle
sparse objects without introducing occlusion errors.

\section{Importance of the Issue}
% Write 2–3 paragraphs that:
% - Explain why solving this problem matters.
% - Distinguish between theoretical significance (advancing knowledge) and practical significance (real-world impact).
% - If relevant, connect to broader societal, educational, or industry needs.

Until now, illustrators have often relied on high sampling to
visually resolve occlusion errors arising from intersecting
objects. Furthermore, to my knowledge, no software exists that
explicitly sorts triangles correctly, let alone lower-dimensional
affine simplices. Most existing approaches either produce
improper sorting or rely on black-box methods. Solving this
problem without using a \(z\)-buffer or ray tracing constitutes
a significant contribution to both the theoretical and practical
aspects of computer graphics, particularly in the domain of
mathematical illustration. Based on years of studying \(3\)D
graphics, I can assert that this algorithm addresses a
fundamental challenge for the future of the field.

From a theoretical standpoint, the method can illustrate
3D scenes composed of points, line segments, and triangles,
demonstrating that a scene of triangles can be partitioned
so as to allow a partial transitive ordering, resulting in
correct visual occlusion.

This work represents a major theoretical advance in the
conceptualization of 3D illustrations. Its synthesis of
elementary linear algebra with computational methods is
unique and represents a significant step beyond the current
state of the art.

\section{Current Limitations and Gaps}
% Write 2–3 paragraphs that:
% - Identify what previous approaches have failed to address (with brief examples).
% - Explain *why* these gaps exist (technical, conceptual, or practical reasons).
% - Show how the gaps lead directly to the core problem you have stated.

Previous approaches largely neglect intersecting simplices and
often produce occlusion errors. In fact, many software packages
do not even triangulate surfaces correctly. For example, drawing
a Klein bottle in Mathematica, Asymptote, or pgfplots results in
intersections that do not occlude properly. My approach resolves
this issue systematically.

Not only does it eliminate intersections, but it also applies
rigorous affine linear algebra for correct occlusion sorting.
These gaps in existing methods persist largely because no one
has previously pursued a comprehensive study of 3D mathematical
illustration with this level of coherence. While many practitioners
achieve impressive results, none combine triangulation, clipping,
and occlusion sorting in a fully integrated framework. Unlike
approaches that rely on splines or high sampling, my package
ensures accurate visualization of sparse geometries.

High sampling is computationally expensive and often insufficient,
since the human eye can detect very small errors. Achieving
the necessary resolution via sampling alone is generally
impractical. My software addresses this challenge directly by
first clipping simplices and then performing occlusion sorting,
ensuring correct visual results efficiently.

\section{Objectives and Goals}
% Write 1 paragraph that:
% - States the main aim of your work.
% - Lists 2–4 specific objectives that break down that aim into achievable outcomes.
% - Use action verbs (develop, evaluate, demonstrate, validate) instead of vague ones (understand, explore).

The objective of this documentation is to present the software I
have developed, demonstrate its correctness, and provide guidance
on its use. Detailed explanations are provided for each of the
novel algorithms, accompanied by examples illustrating how to
apply the software in practice. In addition, practical use cases
are discussed, and the software is validated through tests.

% \section{Research Questions or Hypotheses}
% (Optional — depends on your discipline.)
% - If you use research questions: write them as precise, answerable questions that flow from your objectives.
% - If you use hypotheses: phrase them as testable claims or expectations about outcomes.
% - Keep the list short (usually 2–3 major questions or hypotheses).

\section{Intended Audience}
% Write 1 paragraph that:
% - Identifies who benefits most (students, educators, researchers, industry professionals, etc.).
% - Clarifies how different groups might use the results differently.
% - Helps frame the later discussion of impact and significance.

This documentation is intended for \(3\)D mathematical illustrators
who are already familiar with basic programming, linear algebra,
and Ti\textit{k}Z. While the package primarily supports \(3\)D
mathematical illustrations in \LaTeX{}, its underlying algorithms
have implications beyond the \LaTeX{} community. Computer graphics
researchers will benefit from the first coherent synthesis of a
transitive partial-order algorithm for occlusion sorting of affine
simplices.

Mathematics textbook authors and illustrators will also benefit
from a robust backend for occlusion handling and clipping, without
relying on black-box software that produces incorrect results.
Although some tools generate visually plausible output for users
with limited 3D intuition, subtle errors often remain undetected.
The approach presented here provides a rigorous and systematic
solution, and it is likely that the software will have applications
beyond these immediate use cases.

\section{Summary of the Problem Statement}
% Write 1 short paragraph that:
% - Recaps the central problem in 2–3 sentences.
% - Highlights why it is important and what you aim to achieve.
% - Provides a transition sentence that leads naturally into the Literature Review chapter.

The central problem addressed by this documentation is the
clipping and occlusion ordering of affine simplices. In other
words, it concerns the visualization of \(3\)D parametric scenes.
This software is significant because it represents the first
coherent synthesis for illustrating such scenes.
